% ---------------------------------------------------------------------------
% Glossaire et liste des sigles
% ---------------------------------------------------------------------------

\frontchapter{Glossaire et liste des sigles}

\renewcommand{\arraystretch}{1.25}

% =============================================================================
% LISTE DES SIGLES
% =============================================================================

\section*{Liste des sigles}
\phantomsection
\addcontentsline{toc}{section}{Liste des sigles}

\small

\begin{longtable}{
    @{}
    >{\bfseries}p{0.18\textwidth}
    p{0.76\textwidth}
    @{}
}

\toprule
Sigle & Signification \\
\midrule
\endfirsthead

\toprule
Sigle & Signification \\
\midrule
\endhead

\bottomrule
\endlastfoot

API
&
\textit{Application Programming Interface}. Interface permettant à plusieurs
composants logiciels ou applications de communiquer entre eux.
\\

BDD
&
Base de données.
\\

CLI
&
\textit{Command-Line Interface}. Interface en ligne de commande permettant
d'exécuter et de paramétrer un programme depuis un terminal.
\\

CRS
&
\textit{Coordinate Reference System}. Système de référence de coordonnées
utilisé pour localiser des objets géographiques.
\\

DLL
&
\textit{Dynamic Link Library}. Bibliothèque dynamique utilisée par les
applications Windows pour partager des fonctions et des dépendances.
\\

EPSG
&
Référentiel de codes permettant d'identifier les systèmes de coordonnées et
les transformations géodésiques.
\\

ETL
&
\textit{Extract, Transform, Load}. Processus d'extraction, de transformation et
de chargement de données.
\\

FME
&
\textit{Feature Manipulation Engine}. Plateforme d'intégration et de
transformation de données développée par Safe Software.
\\

GDAL
&
\textit{Geospatial Data Abstraction Library}. Bibliothèque libre permettant de
lire, écrire et transformer de nombreux formats de données géospatiales.
\\

OGR
&
Composant de GDAL principalement consacré à la lecture, à l'écriture et au
traitement des données vectorielles.
\\

RSE
&
Responsabilité sociétale des entreprises.
\\

SDE
&
\textit{Spatial Database Engine}. Technologie Esri utilisée pour accéder à des
données géographiques stockées dans un système de gestion de base de données.
\\

SGBD
&
Système de gestion de base de données.
\\

SIG
&
Système d'information géographique.
\\

SQL
&
\textit{Structured Query Language}. Langage utilisé pour interroger et
manipuler des bases de données relationnelles.
\\

WKT
&
\textit{Well-Known Text}. Représentation textuelle normalisée d'une géométrie
ou d'un système de coordonnées.
\\

\end{longtable}

\normalsize

\vspace{0.6cm}

% =============================================================================
% GLOSSAIRE
% =============================================================================

\section*{Glossaire}
\phantomsection
\addcontentsline{toc}{section}{Glossaire}

\small

\begin{longtable}{
    @{}
    >{\bfseries}p{0.24\textwidth}
    p{0.70\textwidth}
    @{}
}

\toprule
Terme & Définition \\
\midrule
\endfirsthead

\toprule
Terme & Définition \\
\midrule
\endhead

\bottomrule
\endlastfoot

ArcPy
&
Bibliothèque Python fournie par Esri permettant d'automatiser des traitements
réalisés avec ArcGIS et d'accéder aux propriétés des données géographiques
Esri.
\\

Base de données spatiale
&
Base de données capable de stocker, d'interroger et de manipuler des objets
géographiques tels que des points, des lignes et des polygones.
\\

Connecteur
&
Composant logiciel chargé d'établir la communication avec une source de données
et d'en extraire les informations attendues par l'application.
\\

Donnée raster
&
Donnée géographique organisée sous la forme d'une grille de cellules ou de
pixels, utilisée notamment pour représenter des images aériennes, des modèles
numériques ou des cartes continues.
\\

Donnée vectorielle
&
Donnée géographique représentée par des géométries telles que des points, des
lignes ou des polygones, auxquelles peuvent être associés des attributs.
\\

Driver
&
Composant logiciel permettant à GDAL/OGR de lire ou d'écrire un format de
fichier ou une source de données particulière.
\\

Emprise spatiale
&
Rectangle minimal décrivant les limites géographiques d'un jeu de données. Elle
est généralement définie par ses coordonnées minimales et maximales.
\\

Enveloppe convexe
&
Plus petit polygone convexe contenant l'ensemble des géométries d'un jeu de
données.
\\

Géodatabase
&
Structure de stockage de données géographiques utilisée notamment par les
solutions Esri. Elle peut être locale ou hébergée dans un système de gestion de
base de données.
\\

Géomatique
&
Discipline regroupant les méthodes et technologies utilisées pour acquérir,
stocker, traiter, analyser et représenter des données localisées.
\\

Métadonnée
&
Information décrivant une ressource de données, par exemple son nom, sa
structure, son format, sa projection, son emprise ou sa date de modification.
\\

Packaging
&
Processus consistant à regrouper une application et ses dépendances afin de
permettre son installation ou son exécution dans un environnement cible.
\\

PostGIS
&
Extension spatiale du système de gestion de base de données PostgreSQL. Elle
permet de stocker et d'interroger des géométries ainsi que d'effectuer des
traitements spatiaux.
\\

PyInstaller
&
Outil permettant de transformer une application Python et ses dépendances en
un exécutable pouvant être lancé sans installation préalable de Python.
\\

Scan
&
Produit d'ISOGEO chargé d'explorer des sources de données, d'identifier les
ressources qu'elles contiennent et d'extraire les informations techniques
nécessaires à leur documentation dans l'écosystème ISOGEO.
\\

Signature
&
Ensemble d'informations calculées à partir d'une ressource afin de la
caractériser, de la comparer ou de détecter une éventuelle évolution de son
contenu.
\\

Source de données
&
Emplacement ou système à partir duquel Scan analyse des ressources. Une source
peut notamment correspondre à une base de données, une géodatabase, un dossier
ou un fichier géographique.
\\

Test de régression
&
Test visant à vérifier qu'une modification du logiciel n'altère pas un
comportement précédemment validé. Dans ce projet, les résultats Python sont
comparés aux résultats produits par les traitements FME de référence.
\\

Workbench FME
&
Fichier de travail FME décrivant graphiquement une chaîne de lecture, de
transformation et d'écriture de données.
\\

\end{longtable}

\normalsize